% Part: sets-functions-relations
% Chapter: arithmetization
% Section: rationals

\documentclass[../../../include/open-logic-section]{subfiles}

\begin{document}

\olfileid[pt]{sfr}{arith}{rat}
\olsection{De $\Int$ para $\Rat$}

Acabamos de ver como construir os inteiros a partir dos naturais, usando
alguma teoria ingênua dos conjuntos. Veremos agora como construir os racionais
a partir dos inteiros, de modo muito semelhante. As constatações iniciais são:
\begin{enumerate}
	\item Todo racional pode escrever-se na forma $\nicefrac{i}{j}$,
	onde $i$ e $j$ são inteiros, mas $j$ é diferente de zero.
	\item A informação codificada numa expressão $\nicefrac{i}{j}$
	pode igualmente codificar-se num par ordenado $\tuple{ i, j}$.
\end{enumerate}
A abordagem óbvia seria pensar nos racionais \emph{como} pares ordenados de
$\Int \times (\Int \setminus \{0_\Int\})$. Como antes, porém, isso seria um
pouco demasiado ingênuo, pois queremos $\nicefrac{3}{2} = \nicefrac{6}{4}$, mas
$\tuple{ 3, 2}\neq \tuple{ 6,
4}$. Mais geralmente, quereremos o seguinte:
\[
	\nicefrac{a}{b} = \nicefrac{c}{d} \text{ sse } a \times d = b \times c
\]
Para obter isto, definimos uma \emph{relação de equivalência} em $\Int \times
(\Int \setminus \{0_\Int\})$ assim:
\[
	\tuple{ a, b }\Ratequiv \tuple{ c, d} \text{ sse }a \times d = b \times c
\]
Há que verificar que isto é uma relação de equivalência. É muito parecido com o
caso de $\Intequiv$, e deixamo-lo como exercício.
\begin{prob}
Mostre que $\Ratequiv$ é uma relação de equivalência.
\end{prob}
Mas isto permite-nos dizer o seguinte:
\begin{defn}
Os racionais são as classes de equivalência, sob $\Ratequiv$, dos pares de
inteiros (cujo segundo elemento é não nulo). Isto é, $\Rat =
\equivclass{(\Int \times (\Int\setminus \{0_\Int\}))}{\Ratequiv}$.
\end{defn}

Tal como para os inteiros, queremos também definir algumas operações básicas.
Sendo $\equivrep{i,j}{\Ratequiv}$ a classe de equivalência sob $\Ratequiv$ com
$\tuple{i, j}$ como !!a{element}, definimos:
\begin{align*}
	\equivrep{a, b}{\Ratequiv} + \equivrep{c, d}{\Ratequiv} &= \equivrep{ad + bc,  bd}{\Ratequiv}\\
	\equivrep{a, b}{\Ratequiv} \times \equivrep{c, d}{\Ratequiv} &= \equivrep{a  c, b d}{\Ratequiv}.
\intertext{Para definir $r \leq s$ nestes racionais, usamos o fato de que
$r \le s$ sse $s - r$ não é negativo, O.s.\@, $r - s$ pode escrever-se como
$\nicefrac{i}{j}$ com $i$ não negativo e $j$ positivo:}
	\equivrep{a, b}{\Ratequiv} \leq \equivrep{c, d}{\Ratequiv} &\text{ sse }
	\equivrep{c, d}{\Ratequiv} - \equivrep{a, b}{\Ratequiv} =
	\equivrep{i_\Int, j_\Int}{\Ratequiv}
\end{align*}
para algum $i \in \Nat$ e $0 \neq j \in \Nat$.

Depois há que verificar que estas definições se comportam como \emph{devem};
relegamo-lo a \olref[check]{sec}. Mas de fato comportam-se!{} Por fim,
queremos alguma forma de tratar inteiros \emph{como} racionais; logo, para cada
$i \in \Int$, estipulamos $i_\Rat =
\equivrep{i, 1_\Int}{\Ratequiv}$. De novo, verificamos na \olref[check]{sec} que
tudo isto se comporta corretamente.

\begin{prob}
Mostre que $(i + j)_\Rat = i_\Rat + j_\Rat$ e $(i \times j)_\Rat =
i_\Rat \times j_\Rat$ e $i \leq j \liff i_\Rat \leq j_\Rat$, para quaisquer
$i, j \in \Int$.
\end{prob}

\end{document}
